\section{Results}

\subsection{V0 agreement}
Lane A identifies the responsible layer as representation-regime characterization and selects regime-predicate characterization as primary, with the support-two closure test as a complementary probe. Its answer is bound to 13 host capability receipts and the verified evidence items admitted by the frozen protocol.

Lane B identifies the corresponding typed responsibility state, defeats the current-search, donor-family, and method-language alternatives, and selects `COMPUTE:REGIME_CHARACTERIZATION`. Its native decision is bound to the frozen manifest and transition receipts. The revision gate reports the representation split as unresolved because the required predicate obligation has not yet been satisfied.

Under the protocol's compatibility rules the contemporaneous verdict is `AGREE`: the diagnosis and primary move match at the scored level, and the evidence sets overlap by construction while remaining receipt-bound.

\subsection{Deferred coordinate}
After V0, the independently launched frontier lanes produce an exact regime-membership predicate on their stated domains and restore family closure at support two, the two directions prioritized by Lane A and structurally licensed by Lane B. Coordinate 4 is therefore recorded `ALIGNED`. This retrospective score is descriptive; it does not turn the subsequent result into a gold label for all frontier decisions.

\subsection{Historical defects are repaired}
The V1 manuscript documented two real harness defects: malformed successful reasoner content could escape as an unstructured crash, and the deterministic request identity could then be pinned by an immutable semantically invalid receipt. Current regression tests now require structured `HOST_CAPABILITY_FAILED` behavior and provide an explicit, reason-bearing invalid-content archival path that preserves the original bytes before allowing a corrected receipt. These repairs update the implementation status but do not add benchmark instances.

\subsection{Evidence-series gate}
No V1/V2/V3 benchmark series is present on current main. The result set therefore remains exactly one scored question. We report no agreement percentage, confidence interval, calibration measure, or predictive-validity statistic.